\FloatBarrier
\section{本章小结}
\label{sec:direct-methods-summary}

本章回答了章首的问题：求解目标确定以后，怎样不求逆、在有限步内把选定的解算出来。共同的做法是把原问题化为容易求解的问题，并把分解与右端的处理分开，使同一次分解可以服务于多个右端。

对方阵的精确解，消元把求解化为前代与回代。一般非奇异方阵用带部分选主元的LU分解$\bm P\bm A=\bm L\bm U$，乘子按位置存入$\bm L$，后续换行时已存乘子要随行交换；部分选主元对任何非奇异矩阵都能进行，不选主元则要求各顺序主子式非零。

利用对称性可以减少一半的工作。对称正定矩阵用Cholesky分解$\bm A=\bm L\bm L^\top$，每一列的开方量是一个Schur补，正定性恰好保证它为正；分解存在且唯一当且仅当矩阵对称正定，运算量是LU的一半。对称不定矩阵用分块LDL分解$\bm P\bm A\bm P^\top=\bm L\bm D\bm L^\top$，负主元保留在$\bm D$中，零对角元时改用二阶块，使可逆的矩阵总能继续消元。

块结构和低秩修改让这些分解得以复用。分块消元先用被消块$\bm B$的因子求出$\bm Z$与$\bm z$，再解Schur补方程$\bm S\bm v=\bm h$；$\bm B$可逆时，原矩阵非奇异当且仅当$\bm S$非奇异，对称时正定性也传给$\bm S$。低秩更新把修改的影响集中到$k$维辅助量上，只需旧因子上的$k+1$次求解和一个$k$阶方程；更新后的矩阵可逆当且仅当辅助矩阵$\bm K$可逆，Woodbury恒等式和Sherman--Morrison公式是这一等价关系的矩阵形式。两者的节省都有条件：被消块易解，或者已有分解且修改秩很小。

对最小二乘与最小范数问题，三条路线的区别在于怎样处理$\bm A$。正规方程法利用残差正交化为对称正定方程，可以按行累计数据，但使条件数平方，只适合条件较好的满列秩问题。QR求解用正交变换保持残差长度，把目标分成可消去的部分和不可拟合的部分，满列秩时回代$\bm R\hat{\bm x}=\bm c_1$，满行秩时对$\bm A^\top$分解得到最小范数解。

三种QR实现的差别在于怎样计算正交因子。修正Gram--Schmidt直接算出正交基，有限精度下可能失去正交性；Householder反射一次消去整列，选取同号符号避免相消，是稠密问题的标准做法；Givens旋转一次消去一个元素，适合稀疏结构和逐行更新。秩亏或秩未知时，SVD求解在奇异坐标下把目标分离成独立的标量问题，对任意形状和秩给出$\bm A^\dagger\bm b$；加权问题经白化化为普通情形，齐次问题取完整右奇异基中伸缩量最小的方向。

无论用哪种方法，结果都要回到原始数据上按目标检验：精确解看残差是否为零，最小二乘解看残差是否与列空间正交，最小范数解还要看是否与零空间正交。这些检验说明算出的$\bm x$符合目标，却不说明它离真解有多近。主元增长、条件数平方和小奇异值的放大在本章只作了算法层面的说明，结果的可信程度、数值秩的判定以及正则化留给第~\ref{chap:solution-error-regularization}章。矩阵太大而无法分解，或只能通过矩阵与向量的乘积访问矩阵时，就需要从第~\ref{chap:stationary-iterations}章开始的迭代法。
